\documentclass[11pt,a4paper]{article}

\usepackage{geometry}
\geometry{margin=2.5cm}

\usepackage{hyperref}
\usepackage{enumitem}
\usepackage{booktabs}
\usepackage{longtable}
\usepackage{array}
\usepackage{amsmath}

\title{Interface for Observed Passenger Flow Data\\
for Schedule-Based Assignment Models}
\author{Michel Bierlaire \and Evangelos Paschalidis}
\date{April 23, 2026}

\begin{document}
\maketitle

\section{Purpose}

This document specifies a \textbf{standard interface} for providing observed
passenger flow data to be used in the estimation and validation of a
schedule-based public transport assignment model.

The objective is to allow analysts to provide observed flows in a way that is:
\begin{itemize}
    \item consistent with the \textbf{scenario representation} (stops, trips, timetable),
    \item independent of \textbf{internal assignment indexing conventions},
    \item flexible enough to support different types of observations,
    \item robust to \textbf{missing data} and heterogeneous measurement systems,
    \item easily extendable to likelihood-based estimation.
\end{itemize}

This interface describes only the \emph{data format}.  
The mapping between observations and model outputs, as well as the likelihood
specification, will be handled by separate components.

\section{Design principles}

\subsection{Separation from assignment internals}

The assignment algorithm operates on time-expanded graphs and internal link
indices. These indices must remain \emph{internal} to the model implementation.
Analysts providing data should not need to know them.

Therefore, observations are expressed exclusively in terms of:
\begin{itemize}
    \item scenario-level identifiers (stop, line, trip),
    \item timetable times,
    \item measurement metadata.
\end{itemize}

\subsection{Event-based timing}

Time bins in the scenario represent desired departure times and are not tied
to timetable events. Observations of passenger flows, however, are associated
with \textbf{events in the timetable} (arrivals, departures, vehicle loads).

Hence:
\begin{itemize}
    \item All observed flows are linked to timetable events.
    \item Time is expressed as clock time (\texttt{HH:MM:SS}).
    \item Observations refer to time windows on the service day.
\end{itemize}

\subsection{Partial and heterogeneous data}

Not every link in the network has data. Observations may include:
\begin{itemize}
    \item vehicle load counts,
    \item boarding counts,
    \item alighting counts,
    \item aggregated measurements.
\end{itemize}

The interface must therefore:
\begin{itemize}
    \item allow missing data naturally,
    \item support multiple measurement methods,
    \item remain simple for analysts to fill in.
\end{itemize}

\subsection{Separation of data and statistical assumptions}

The data file should contain \emph{observations only}.  
It should not include statistical parameters such as variances or weights.

Instead, each observation is tagged with a \texttt{method\_id}.  
The likelihood wrapper will interpret this identifier and assign the appropriate
measurement model.

\section{File format}

Observed flows are provided in a table file:
\begin{itemize}
    \item \texttt{observations.csv} (recommended), or
    \item Parquet with identical schema.
\end{itemize}

Each row represents one measurement.

\section{Columns}

\subsection{Identification and metadata}

\begin{longtable}{p{4cm}p{3cm}p{8cm}}
\toprule
Column & Type & Description \\
\midrule
\texttt{obs\_id} & string & Unique identifier of the observation. \\
\texttt{method\_id} & string & Identifier of measurement method or sensor. \\
\texttt{obs\_type} & string & Type of observation. Currently: \texttt{instant\_load}. \\
\texttt{value} & float & Observed passenger count. Must be $\ge 0$. \\
\bottomrule
\end{longtable}

\subsection{Time window}

\begin{longtable}{p{4cm}p{3cm}p{8cm}}
\toprule
Column & Type & Description \\
\midrule
\texttt{time\_start} & HH:MM:SS & Start of observation window. \\
\texttt{time\_end} & HH:MM:SS & End of observation window (exclusive). \\
\bottomrule
\end{longtable}

Interpretation: the measurement refers to timetable events within
$[\texttt{time\_start},\texttt{time\_end})$.

If an observation is instantaneous, a short window (e.g.\ 30--60s) should be used.

\subsection{Location and event}

Each observation must identify where in the timetable it applies.

\subsubsection*{Trip-based mode (preferred)}

\begin{longtable}{p{4cm}p{3cm}p{8cm}}
\toprule
Column & Type & Description \\
\midrule
\texttt{trip\_id} & string & Trip identifier from the timetable. \\
\texttt{stop\_id} & string & Stop identifier. \\
\texttt{event} & string & \texttt{ARR} or \texttt{DEP}. \\
\bottomrule
\end{longtable}

\subsubsection*{Line-based mode (aggregate)}

Used when trip identity is unknown.

\begin{longtable}{p{4cm}p{3cm}p{8cm}}
\toprule
Column & Type & Description \\
\midrule
\texttt{line\_id} & string & Line identifier. \\
\texttt{stop\_id} & string & Stop identifier. \\
\texttt{event} & string & \texttt{ARR} or \texttt{DEP}. \\
\bottomrule
\end{longtable}

\textbf{Rule:} exactly one of \texttt{trip\_id} or \texttt{line\_id} must be provided.

\subsection{Optional metadata}

\begin{longtable}{p{4cm}p{3cm}p{8cm}}
\toprule
Column & Type & Description \\
\midrule
\texttt{campaign\_id} & string & Measurement campaign. \\
\texttt{source} & string & Sensor or file source. \\
\texttt{notes} & string & Free-text notes. \\
\texttt{stop\_sequence} & int & Optional disambiguation if a trip visits a stop multiple times. \\
\bottomrule
\end{longtable}

\section{Instantaneous load semantics}

For \texttt{obs\_type = instant\_load}, the value represents observed onboard
passengers near a timetable event.

The exact interpretation (e.g.\ load after departure vs.\ before arrival) is
defined by the measurement method and handled by the mapping layer.

\section{Validation rules}

The loader should verify:

\begin{itemize}
    \item \texttt{obs\_id} unique.
    \item \texttt{value} $\ge 0$.
    \item Valid time format \texttt{HH:MM:SS}.
    \item \texttt{time\_end} $>$ \texttt{time\_start}.
    \item \texttt{stop\_id} exists in scenario.
    \item \texttt{trip\_id} or \texttt{line\_id} exists.
    \item Exactly one of them is provided.
    \item \texttt{event} in \{\texttt{ARR}, \texttt{DEP}\}.
\end{itemize}

\section{Example}

\begin{verbatim}
obs_id,method_id,obs_type,value,time_start,time_end,trip_id,line_id,stop_id,event
L001,APC_V1,instant_load,37,08:00:00,08:02:00,T12,,A,DEP
L002,APC_V1,instant_load,22,08:10:00,08:12:00,T12,,B,ARR
L101,MANUAL_V1,instant_load,18,17:30:00,17:45:00,,L1,C,DEP
\end{verbatim}

\section{Future extensions}

The interface is designed to support:
\begin{itemize}
    \item boarding counts,
    \item alighting counts,
    \item link flows,
    \item likelihood-based estimation.
\end{itemize}

Mapping from observations to model quantities will be handled by
dedicated mapper components, using:
\begin{itemize}
    \item the scenario,
    \item the assignment outputs,
    \item the \texttt{method\_id}.
\end{itemize}



\end{document}
